\section{结论}\label{sec:conclusion}
本文以最小二乘消去未知宽带波形，将同步单天线站的条件搜索、残差发现和联合精修统一到同一位置目标。多候选联合AP保留粗搜索中尚未占优的候选，以精修后的拟合残差择优；波形消元块NOMP保留残差发现和循环更新，通过消元曲率减少位置--信号系数耦合造成的重复小步。

已完成的配对实验表明，候选保留能够恢复单候选的错误区域，联合精修在相同精度要求下减少了目标与梯度调用。受控实验进一步将固定系数更新的慢收敛与曲率机制联系起来，支持以消元目标改善局部更新。两种方法分别从候选可靠性和更新效率出发，为有限观测下的多源直接定位提供了可解释的求解策略。扩大布局的统一比较将用于确定最终性能范围与推荐配置；未知源数及更复杂传播可在这一基础上继续研究。
